\section{Grader 体系：把产品约束翻译为训练信号}

讲者把 grader 称为核心工程资产。因为在 RL 训练中，模型是否变好取决于 ``奖励是否真正代表期望行为''。如果 grader 只看 pass rate，就会错过用户体验和产品规范。

课程中提到的常见 grader 包括：unit test verifier、LLM-as-judge、rubric-based grading、格式约束、长度约束、时延约束、行为合规约束等。它们通常并非单一评分器，而是一个多维评分栈。

\begin{importantbox}{``定义好 grader = 解决一半问题''}
这是本讲最具工程味的结论之一。因为一旦评分目标不对，后续最优策略只会把系统推向错误方向，形成高分低能。
\end{importantbox}

\begin{table}[H]
\centering
\begin{tabular}{p{2.5cm}p{3.2cm}p{4.8cm}}
\toprule
\textbf{Grader 类型} & \textbf{适用场景} & \textbf{设计注意点} \\
\midrule
Verifier（程序可检验） & 单测、编译、格式检查 & 防止模型篡改测试环境 \\
Model Judge & 文本质量、解释充分性 & 需校准偏差，避免同模偏好 \\
Rubric Scoring & 多维行为约束 & 维度拆分清晰，权重可解释 \\
Constraint Grader & 时延、长度、成本 & 与质量指标联合优化而非独立优化 \\
\bottomrule
\end{tabular}
\caption{课堂提及的 grader 组件化思路}
\end{table}

讲者提到 grader 增多后会出现冲突与依赖。例如某个格式 grader 可能与长度控制冲突，某些安全 grader 必须在功能 grader 之后执行。工程上需要显式建模依赖图，而不是简单线性叠加。

\begin{warningbox}{评分器冲突是常态，不是异常}
多目标任务里，``每加一个 grader 都增益'' 是理想化假设。真实系统常见情形是新增约束导致旧指标回退，需要重新分层权重和训练阶段。
\end{warningbox}

\begin{knowledgebox}{讲者隐含的三步法}
\begin{enumerate}
    \item 从真实用户反馈中抽取失败模式，确定 ``要评什么''；
    \item 为每类失败模式设计可执行 grader；
    \item 监控 grader 之间冲突，按训练阶段动态启用。
\end{enumerate}
\end{knowledgebox}

\subsection{本章小结}
在 agent 训练中，grader 体系等价于产品规范的形式化表达。没有高质量 grader，就没有可持续的 RL 改进。

